\chapter{چارچوبِ وب و سامانهٔ چیدمان}
\label{ch:web}

سلام تنها برای برنامه‌های متنی نیست؛ یک \textbf{سامانهٔ چیدمان} (layout) دارد که
با آن می‌توانید صفحه‌های وب بسازید آن هم به زبانِ فارسی و راست‌به‌چپ! در این
فصل یاد می‌گیریم چگونه با سلام یک صفحهٔ وبِ کامل و زیبا طراحی کنیم.

\section{نخستین صفحهٔ وب}

برنامه‌های چیدمان به‌جای \code{تابع اصلی}، با واژهٔ \code{صفحه} (معادلِ
\code{layout}) آغاز می‌شوند. این مثالِ واقعی یک صفحهٔ ساده می‌سازد:

\begin{salam}
صفحه:
    عنوان = "سلام دنیا"
    جهت = "راست به چپ"
    جعبه:
        محتوا = "درود"
        رنگ = "آبی"
        پس‌زمینه = "سفید"
    پایان
پایان
\end{salam}

برای ساختنِ صفحهٔ وب از این کد، از فرمانِ \code{layout build} استفاده می‌کنیم:

\begin{salamen}[language={}]
salam layout build page.salam
\end{salamen}

این فرمان سه فایل می‌سازد: \code{page.html}، \code{page.css} و \code{page.js}.
کافی است \code{page.html} را در مرورگر باز کنید تا صفحهٔ خود را ببینید! اگر یک
فایلِ یکپارچه می‌خواهید، گزینهٔ \code{--inline} را بیفزایید:

\begin{salamen}[language={}]
salam layout build page.salam --inline
\end{salamen}

\section{عنصرها و ویژگی‌ها}

هر صفحه از \textbf{عنصر}ها ساخته می‌شود (مانندِ \code{جعبه}، \code{سرصفحه}،
\code{دکمه})، و هر عنصر \textbf{ویژگی}هایی دارد (مانندِ \code{رنگ}،
\code{محتوا}، \code{پس‌زمینه}). الگوی همیشگی چنین است:

\begin{center}
\textit{نامِ‌عنصر}\code{:} \quad \textit{ویژگی} \code{=} \textit{مقدار} \quad \ldots \quad \code{پایان}
\end{center}

این مثالِ واقعی که در پوشهٔ مثال‌های سلام هست یک صفحهٔ کامل‌تر با سرصفحه
می‌سازد:

\begin{salam}
صفحه:
    عنوان = "برنامه‌ی من"
    زبان = "فارسی"
    جهت = "راست به چپ"
    سرصفحه:
        پس‌زمینه = "آبی"
        رنگ = "سفید"
        فاصله داخلی = "20px"
        سرتیتر: اندازه = 1 محتوا = "خوش آمدید" اندازه_قلم = "32px" پایان
    پایان
    جعبه:
        محتوا = "این یک صفحه‌ی نمونه به زبان فارسی است."
        رنگ = "سیاه"
    پایان
پایان
\end{salam}

می‌بینید که عنصرها می‌توانند \emph{تودرتو} باشند: \code{سرتیتر} درونِ
\code{سرصفحه} و آن درونِ \code{صفحه}. این درست مانندِ ساختارِ یک صفحهٔ وبِ واقعی
است که بخش‌هایی درونِ بخش‌های دیگر دارد.

\begin{note}
ویژگی‌های \code{زبان} و \code{جهت} به‌ویژه برای فارسی مهم‌اند. با
\code{جهت = "راست به چپ"} کلِ صفحه راست‌به‌چپ چیده می‌شود همان‌طور که متنِ
فارسی باید باشد. این یکی از جاهایی است که پشتیبانیِ بومیِ سلام از فارسی
می‌درخشد.
\end{note}

\section{عنصرهای پرکاربرد}

سلام عنصرهای فراوانی برای ساختنِ صفحه دارد. مهم‌ترین‌ها:

\begin{center}
\begin{tabular}{l l l}
\toprule
\textbf{فارسی} & \textbf{معادلِ HTML} & \textbf{کاربرد} \\
\midrule
\code{صفحه} & \lr{html/body} & ریشهٔ صفحه \\
\code{سرصفحه} & \lr{header} & بخشِ بالای صفحه \\
\code{پاصفحه} & \lr{footer} & بخشِ پایینِ صفحه \\
\code{سرتیتر} & \lr{h1..h6} & تیتر \\
\code{جعبه} & \lr{div} & یک بخشِ کلی \\
\code{بند} & \lr{p} & پاراگراف \\
\code{فهرست} & \lr{ul} & فهرست \\
\code{مورد} & \lr{li} & یک عضوِ فهرست \\
\code{دکمه} & \lr{button} & دکمه \\
\code{فرم} & \lr{form} & فرمِ ورودی \\
\code{ورودی} & \lr{input} & جعبهٔ ورودی \\
\code{پیوند} & \lr{a} & پیوند \\
\code{تصویر} & \lr{img} & تصویر \\
\bottomrule
\end{tabular}
\end{center}

\subsection{فهرست}

این مثالِ واقعی یک فهرست از میوه‌ها می‌سازد:

\begin{salam}
صفحه:
    فهرست:
        مورد: محتوا = "سیب" پایان
        مورد: محتوا = "پرتقال" پایان
        مورد: محتوا = "موز" پایان
    پایان
پایان
\end{salam}

\subsection{دکمه}

\begin{salam}
صفحه:
    دکمه: نوع = "دکمه" محتوا = "کلیک کنید" پس‌زمینه = "آبی" رنگ = "سفید" پایان
پایان
\end{salam}

\section{ساختنِ یک فرم}

فرم‌ها به کاربر اجازه می‌دهند داده وارد کند. این مثالِ واقعی یک فرمِ ورود با
جعبهٔ گذرواژه می‌سازد:

\begin{salam}
صفحه:
    فرم:
        ورودی: نوع = "گذرواژه" نام = "pass" جانگهدار = "رمز عبور" پایان
        دکمه: نوع = "ارسال" محتوا = "ورود" پایان
    پایان
پایان
\end{salam}

عنصرِ \code{ورودی} ویژگی‌های کلیدی دارد:
\begin{itemize}
  \item \code{نوع}: گونهٔ ورودی \code{"متن"}، \code{"گذرواژه"}،
        \code{"ایمیل"}، \code{"ارسال"} و مانندِ این‌ها.
  \item \code{نام}: نامی که داده با آن فرستاده می‌شود.
  \item \code{جانگهدار}: متنِ راهنمایی که پیش از تایپِ کاربر نشان داده می‌شود.
\end{itemize}

\begin{tip}
نوعِ \code{"گذرواژه"} به مرورگر می‌گوید که نویسه‌های تایپ‌شده را به‌جای متن، با
نقطه نشان دهد تا کسی از روی شانهٔ کاربر رمز را نبیند. این یک نکتهٔ کوچک اما
مهمِ امنیتی و تجربهٔ کاربری است.
\end{tip}

\section{سبک‌دهی و ظاهر}

ویژگی‌هایی مانندِ \code{رنگ}، \code{پس‌زمینه}، \code{فاصله\_داخلی} و
\code{اندازه\_قلم} ظاهرِ عنصرها را تعیین می‌کنند. سلام این‌ها را به CSSِ
استاندارد ترجمه می‌کند. مقدارهای رنگ می‌توانند نام‌های فارسی (\code{"آبی"}،
\code{"سفید"}، \code{"سیاه"}) یا کدهای استاندارد باشند.

\begin{tip}
برای ویژگی‌های اندازه‌ای مانندِ \code{فاصله\_داخلی}، \code{حاشیه}،
\code{عرض}، \code{ارتفاع} و \code{اندازه\_قلم} لازم نیست همیشه واحد را
بنویسید: اگر فقط یک عدد بدهید (مثلاً \code{فاصله داخلی = 24})، سلام
به‌طورِ خودکار \code{px} (پیکسل) را می‌افزاید. واحدهای دیگر هم پشتیبانی
می‌شوند، چه با فاصله نوشته شوند چه بی‌فاصله، و چه به‌صورتِ اختصاریِ
لاتین (\code{cm}، \code{mm}، \code{em}، \code{rem}، \code{vh}، \code{\%}
و مانندِ این‌ها) چه با نامِ فارسیِ کامل مانندِ
\code{"3 سانتی متر"} یا \code{"3 سانتی‌متر"} (با نیم‌فاصله) و
\code{"5 میلی متر"}. حتی مقدارهای چندتایی مانندِ \code{"10 20"} به
\code{10px 20px} تبدیل می‌شوند.
\end{tip}

\begin{deep}[نگاهی به درون: از سلام تا HTML]
وقتی \code{salam layout build} را اجرا می‌کنید، سلام درختِ عنصرهای شما را
می‌پیماید و برای هر کدام، HTMLِ متناظر را تولید می‌کند؛ ویژگی‌های ظاهری را به
کلاس‌های CSS تبدیل می‌کند (با نام‌هایی که به‌طورِ خودکار از روی محتوا ساخته
می‌شوند تا یکتا باشند)؛ و رفتارها را در یک فایلِ JavaScript می‌گذارد. نتیجه یک
صفحهٔ وبِ کاملاً استاندارد است که در هر مرورگری کار می‌کند اما شما آن را به
فارسی و با واژگانِ آشنا نوشتید.
\end{deep}

\section{خلاصهٔ فصل}

\begin{itemize}
  \item برنامه‌های چیدمان با \code{صفحه} آغاز می‌شوند، نه \code{تابع اصلی}.
  \item صفحه از عنصرهای تودرتو ساخته می‌شود؛ هر عنصر ویژگی‌هایی دارد.
  \item با \code{salam layout build} صفحه به HTML/CSS/JS تبدیل می‌شود.
  \item ویژگی‌های \code{زبان} و \code{جهت} پشتیبانیِ بومی از فارسیِ راست‌به‌چپ
        می‌دهند.
  \item عنصرهای فراوانی هست: \code{جعبه}، \code{فهرست}، \code{دکمه}،
        \code{فرم}، \code{ورودی} و بیشتر.
  \item ویژگی‌های ظاهری به CSS ترجمه می‌شوند.
\end{itemize}

\section{تمرین‌ها}

\exercise یک صفحهٔ وب بسازید که نام و یک پیامِ خوشامد را نشان دهد، با عنوان و
جهتِ راست‌به‌چپ.

\exercise صفحه‌ای با یک فهرستِ سه‌موردی از کارهای روزانه‌تان بسازید.

\exercise یک فرمِ ثبت‌نام با ورودی‌های نام، ایمیل و گذرواژه و یک دکمهٔ ارسال
طراحی کنید.

\exercise صفحه‌ای با یک سرصفحهٔ رنگی و یک جعبهٔ محتوا بسازید و آن را با
\code{--inline} بسازید.

\exercise یک کارتِ معرفی بسازید: جعبه‌ای با یک سرتیتر (نام)، یک بند (توضیح) و
یک دکمه (تماس).
